\section{真实场景应用案例}\label{sec:case}

\subsection{场景设计：开发助手日常使用（办公自动化）}

选取「基于 OS Agent 的开发 / 办公自动化助手」作为真实落地场景，覆盖三类记忆能力的
完整闭环。偏好沉淀与复用方面，用户在多轮会话中反复于同类任务选择某一工具、要求简洁
输出格式，系统从工具执行结果与用户行为中自动提取工具选择与输出风格偏好，跨会话、
跨场景（终端场景 / 文档场景）持续生效；用户手动修改配置后，偏好经版本化更新并可
回溯。知识沉淀与智能调用方面，一次完整的「构建—排错—修复」流程被沉淀为可复用模板
（含前置条件与步骤），五次真实执行证据且成功率达标并经人工复核后激活，后续同类故障
一键复用；新旧事实（如项目构建参数变更）经冲突处理自动接替，版本史可查。安全合规
方面，记忆过程中出现的手机号、邮箱等被自动脱敏；演示「忘掉我的所有住址，但保留城市
偏好」——排除项保留、目标数据物理删除、磁盘零残留。

\subsection{实测记录}

案例执行记录、关键步骤截图与启用 / 关闭记忆的前后对比 \todo{真实场景案例执行并
记录（步骤截图、前后对比、用户反馈）}。

效果演示视频 \todo{录制偏好捕捉、知识整合、检索复用三段演示视频}。

\subsection{应用价值}

该场景直接对应赛题「办公自动化、智能协作、个性化助手」的应用方向：偏好记忆减少
重复指令与工具选择成本，知识模板把一次性排错经验转化为可复用资产，精准遗忘满足
个人信息保护的合规要求。量化收益（任务完成时间缩短比、指令次数下降比）
\todo{案例量化收益统计}。
